\section{Discussion}
\label{sec:discussion}

The central message of \textsc{TAG} is that the utility of a reliable
instruction--response pair depends on the current model state and need
not be captured by a scalar reward alone. We augment the
difficulty--uncertainty carrier with a state-conditioned representation
anchor: principal directions are extracted from the current probe's
within-sequence representation displacements, and normalized candidate
projections provide a bounded multiplicative modulation. This connects
static representation-geometric selection and dynamic scalar-reward
selection under a single ranking rule with one exposed knob,
\(\lambda\).

The theoretical claim is deliberately local. Under covariance-drift
and eigengap conditions, the Davis--Kahan analysis bounds the
one-refresh change of the sign-oriented principal direction at the
optimiser scale. It does not by itself establish stability of candidate
projections or of the full ranking: candidate activations also change,
and pool-wise normalization can alter relative margins. Separately,
the factor
\(1+\lambda\,\widetilde a_i^{(t)}\in[1,1{+}\lambda]\) limits the
anchor projection's direct modulation of the carrier and makes that
role independently ablatable.
The \(\lambda\)-ablation in Fig.~\ref{fig:lambda-ablation} traces out
this structure empirically: performance peaks at the empirical default
\(\lambda{=}1.0\) and declines under stronger modulation, consistent
with using the anchor projection as an auxiliary signal.

\paragraph{Practical implications.}
The scoring pass computing the difficulty--uncertainty carrier also
produces the hidden states needed for the anchor, leaving only
per-layer probe PCA and candidate--anchor projections as additional
work. This yields the empirical profile of Table~\ref{tab:efficiency}
on the main setting: a $-33\%$ reduction in end-to-end wall-clock cost
relative to Data Agent at a higher average score, without an external
LLM judge, curated seed set, or separate policy-learning loop. For
repeatedly refreshed SFT pipelines that must
remain adaptive but cannot budget API calls or human curation, this is
the practical lever.

\paragraph{Scope and future directions.}
We target a conservative setting---text-only SFT, single-sequence
instruction--response examples, rank-one anchor per layer---that keeps
the anchor interpretable and the stability analysis clean. The
robustness regimes in Tables~\ref{tab:backbone-robustness}--%
\ref{tab:dataset-transfer} show that gains from anchor modulation
narrow when the source pool is already strong, as in Evol-Instruct on
larger backbones; we read this as a natural headroom effect of
high-quality pools rather than an anchor failure, and as a useful
boundary case for future analysis. The broader principle generalises:
multi-turn dialogue suggests turn-aware anchors, multimodal tuning
suggests modality-specific anchors, and multi-skill samples suggest
rank-\(r\) representation subspaces. Extending the construction to
reinforcement fine-tuning is also natural but non-trivial: prompts,
rollouts, and reward-conditioned states define a different anchor space,
and the stability analysis would have to be redone for that setting.

